% ============================================================================
%  Section I — Introduction  (IEEE TIM, IEEEtran)
%  Passive voice / third person throughout. Requires: none beyond base
%  (uses \IEEEPARstart from IEEEtran, and \itemize).
% ============================================================================

\section{Introduction}
\label{sec:intro}

\IEEEPARstart{I}{nduction} motors account for a dominant share of the electrical
energy consumed by industry and form the mechanical backbone of countless
processes; their unexpected failure causes costly downtime and, in critical
applications, safety hazards. Continuous condition monitoring (CM) and early
fault diagnosis are therefore central to modern maintenance
practice~\cite{gangsar2020signal, hamani2025advancements}. Among the available
modalities, vibration analysis is the most established, because the mechanical
and electrical faults of interest---bearing defects, rotor asymmetries, and
supply imbalances---imprint characteristic signatures on the machine's vibration
response~\cite{yang2021fault, zhang2020deep}.

The prevailing paradigm relies on industrial-grade piezoelectric accelerometers
sampled at tens to hundreds of kilohertz, which resolve the high-frequency
bearing defect frequencies exploited by spectral and envelope analysis. This is
reflected in the public benchmarks that anchor the field, acquired at
$8$--$200$~kHz~\cite{lee2024multidomain, sehri2024ottawa, thuan2023hust,
huang2018bearing}. Although accurate, such instrumentation is costly, requires
intrusive installation, and scales poorly to the vast population of small,
low-value motors that individually do not justify a dedicated monitoring channel
yet collectively dominate a plant's motor inventory. Closing this coverage gap
motivates a shift toward inexpensive, wireless, and readily deployable sensing.

Consumer micro-electro-mechanical-systems (MEMS) inertial sensors, epitomised by
those embedded in smartphones, offer such a platform at negligible marginal
cost~\cite{grouios2022accelerometers}. Motor fault detection has been demonstrated
from smartphone- and MEMS-acquired vibration~\cite{naha2017mobile,
ertargin2023deep, ertargin2025automated}, and dedicated low-cost MEMS nodes now
integrate on-board processing with wireless telemetry for machine
diagnostics~\cite{pedotti2020lowcost, jakobsen2024lowcost}. These sensors,
however, operate at sampling rates one to three orders of magnitude below
industrial practice---$100$~Hz in the present case---so that, by the Nyquist
criterion~\cite{nyquist1924certain, shannon1948mathematical}, the high-frequency
signatures underpinning classical bearing diagnosis are not observable. A
fundamental and largely unresolved question follows: what fault information is
genuinely recoverable from such a constrained sensor, and does the apparent
performance transfer to a real deployment?

The second half of that question is easily overlooked. Learning-based fault
diagnosis routinely reports near-perfect accuracy, yet a growing body of evidence
attributes much of this success to \emph{data leakage}---information shared
between training and test partitions that inflates measured accuracy without
improving generalisation~\cite{kapoor2023leakage}. In vibration diagnosis, leakage
arises when overlapping analysis windows from a single continuous recording are
divided across partitions, or when recording-specific characteristics are learned
in place of fault features; enforcing component- or recording-wise separation has
been shown to reduce reported accuracy substantially~\cite{wheat2024impact,
hendriks2022towards}. This scrutiny has so far been directed at industrial-grade,
high-rate data. For low-rate smartphone datasets---in which each operating
condition is captured in one continuous acquisition and heavy window overlap is
common---the risk is arguably greater, yet random-split evaluation remains the
norm, including in the baseline accompanying the dataset examined
here~\cite{ertargin2026smartphone}.

A third consideration is deployment. The Industrial Internet of Things envisions
dense networks of inexpensive nodes that stream, or locally distil, diagnostic
information over constrained wireless links~\cite{yousuf2024iot}, and recent
TinyML implementations have shown that lightweight classifiers can run on
microcontrollers at millisecond latency and sub-millijoule
energy~\cite{gao2025edge, asutkar2023tinyml}. Such work establishes that edge
diagnosis is feasible, but it generally presupposes industrial-bandwidth sensing
and does not characterise the coupled trade-off between diagnostic accuracy and
the sampling, communication, and computational cost that governs an ultra-low-rate
node. Whether, for a smartphone-class sensor, the low sampling rate is a liability
to be tolerated or an asset to be exploited has not been quantified.

This paper addresses these three gaps jointly, using a public smartphone-acquired
vibration dataset of an induction motor recorded across six health states, three
supply frequencies, and two load conditions at
$100$~Hz~\cite{ertargin2026smartphone, ertargin2026dataset}. Rather than propose a
new classifier, the study foregrounds measurement rigour: genuine diagnostic
capability is separated from evaluation optimism, the consumer sensor is
characterised as a measurement instrument, and the accuracy--resource trade-off
that determines an IoT deployment is mapped. The principal contributions are as
follows.

\begin{itemize}
\item A \emph{deployment-realistic evaluation methodology} for windowed low-rate
vibration CM---comprising leakage-free recording-wise and cross-condition
(cross-speed and cross-load) protocols---together with a quantification of the
optimism induced by the prevailing random-split practice, which is shown to
overstate accuracy by approximately $40$ to $45$ percentage points.
\item A \emph{measurement characterisation} of the smartphone MEMS channel that
establishes its effective resolution, noise floor, and usable band, and reveals
that the on-device gravity compensation discards diagnostically useful
low-frequency content, so that the raw channels are preferable for condition
monitoring.
\item An \emph{accuracy--resource trade-off characterisation} that identifies a
resource-frugal operating point outperforming the naive configuration under
honest evaluation, together with a communication link budget and an edge-model
footprint that translate the trade-off into deployable operating points for
Bluetooth Low Energy, Zigbee, NB-IoT, and LoRaWAN nodes.
\item A \emph{reproducible and deployment-honest baseline} for low-cost
IoT-based induction-motor condition monitoring, with all code and configurations
released.
\end{itemize}

The remainder of this paper is organised as follows.
Section~\ref{sec:related} reviews related work on vibration-based condition
monitoring, low-cost MEMS sensing, IoT and edge diagnosis, and evaluation
methodology. Section~\ref{sec:dataset} describes the dataset and measurement
setup, and Section~\ref{sec:charact} characterises the sensor and its signals.
Section~\ref{sec:method} details the processing chain and the evaluation
protocols. Section~\ref{sec:results} reports the results,
Section~\ref{sec:discussion} discusses their implications and limitations, and
Section~\ref{sec:conclusion} concludes.
